\documentclass[aps,prb,twocolumn]{revtex4-2}
\usepackage{amsmath}
\begin{document}
\title{Normal-Mode Splitting in Two Coupled Quartic Oscillators}
\author{A. Researcher}
\begin{abstract}
We compute the first-order frequency shift of two linearly coupled quartic oscillators using canonical perturbation theory.
\end{abstract}
\maketitle
\section{Model Hamiltonian}\label{sec:ham}
Consider two identical oscillators of unit mass with frequency $\omega_0$, linear coupling $g$, and quartic anharmonicity $\lambda$:
\begin{equation}\label{eq:ham-total}
H = \frac{1}{2}(p_1^2 + p_2^2) + \frac{\omega_0^2}{2}(q_1^2 + q_2^2) + g q_1 q_2 + \frac{\lambda}{4}(q_1^4 + q_2^4).
\end{equation}
\subsection{Normal coordinates}\label{sec:normal}
Introducing symmetric and antisymmetric coordinates $Q_{\pm} = (q_1 \pm q_2)/\sqrt{2}$, the quadratic part diagonalizes:
\begin{align}
\omega_+^2 &= \omega_0^2 + g, \label{eq:omega-plus} \\
\omega_-^2 &= \omega_0^2 - g, \label{eq:omega-minus} \\
\delta\omega^2 &= \omega_+^2 - \omega_-^2 = 2g. \nonumber
\end{align}
\section{Perturbative shift and numerical comparison}\label{sec:shift}
To first order in $\lambda$, the averaged quartic potential shifts the symmetric frequency by
\begin{equation}\label{eq:shift}
\Delta\omega_+ = \frac{3\lambda}{8\omega_+} A_+^2.
\end{equation}
\begin{table}[h]
\caption{Comparison of perturbative and numerical frequency shifts for $\omega_0=1$, $g=0.2$.}\label{tab:shifts}
\begin{tabular}{ccc}
$\lambda$ & Perturbative & Numerical \\
0.01 & 0.0034 & 0.0034 \\
\end{tabular}
\end{table}
\begin{figure}[t]
\caption{Splitting $\omega_+ - \omega_-$ versus coupling $g$.}\label{fig:split}
\end{figure}
\bibliography{oscillators}
\end{document}
